\documentclass[10pt,a4paper]{amsart}
\usepackage{amsmath,amssymb,amsthm,graphicx}
% For commutative diagrams we use Xy-pic
\usepackage[all]{xy}

% We use 2cell for 2-commutative diagrams.
\xyoption{2cell}
% Theorem environments.
%
\theoremstyle{plain}
\newtheorem{theorem}[subsection]{Theorem}
\newtheorem{proposition}[subsection]{Proposition}
\newtheorem{lemma}[subsection]{Lemma}

\theoremstyle{definition}
\newtheorem{definition}[subsection]{Definition}
\newtheorem{example}[subsection]{Example}
\newtheorem{exercise}[subsection]{Exercise}
\newtheorem{situation}[subsection]{Situation}

\theoremstyle{remark}
\newtheorem{remark}[subsection]{Remark}
\newtheorem{remarks}[subsection]{Remarks}

\numberwithin{equation}{subsection}

% Macros
%
\def\lim{\mathop{\mathrm{lim}}\nolimits}
\def\colim{\mathop{\mathrm{colim}}\nolimits}
\def\Spec{\mathop{\mathrm{Spec}}}
\def\Hom{\mathop{\mathrm{Hom}}\nolimits}
\def\Ext{\mathop{\mathrm{Ext}}\nolimits}
\def\SheafHom{\mathop{\mathcal{H}\!\mathit{om}}\nolimits}
\def\SheafExt{\mathop{\mathcal{E}\!\mathit{xt}}\nolimits}
\def\Sch{\mathit{Sch}}
\def\Mor{\mathop{\mathrm{Mor}}\nolimits}
\def\Ob{\mathop{\mathrm{Ob}}\nolimits}
\def\Sh{\mathop{\mathit{Sh}}\nolimits}
\def\NL{\mathop{N\!L}\nolimits}
\def\CH{\mathop{\mathrm{CH}}\nolimits}
\def\proetale{{pro\text{-}\acute{e}tale}}
\def\etale{{\acute{e}tale}}
\def\QCoh{\mathit{QCoh}}
\def\Ker{\mathop{\mathrm{Ker}}}
\def\Im{\mathop{\mathrm{Im}}}
\def\Coker{\mathop{\mathrm{Coker}}}
\def\Coim{\mathop{\mathrm{Coim}}}

% Boxtimes
%
\DeclareMathSymbol{\boxtimes}{\mathbin}{AMSa}{"02}

%
% Macros for moduli stacks/spaces
%
\def\QCohstack{\mathcal{QC}\!\mathit{oh}}
\def\Cohstack{\mathcal{C}\!\mathit{oh}}
\def\Spacesstack{\mathcal{S}\!\mathit{paces}}
\def\Quotfunctor{\mathrm{Quot}}
\def\Hilbfunctor{\mathrm{Hilb}}
\def\Curvesstack{\mathcal{C}\!\mathit{urves}}
\def\Polarizedstack{\mathcal{P}\!\mathit{olarized}}
\def\Complexesstack{\mathcal{C}\!\mathit{omplexes}}
% \Pic is the operator that assigns to X its picard group, usage \Pic(X)
% \Picardstack_{X/B} denotes the Picard stack of X over B
% \Picardfunctor_{X/B} denotes the Picard functor of X over B
\def\Pic{\mathop{\mathrm{Pic}}\nolimits}
\def\Picardstack{\mathcal{P}\!\mathit{ic}}
\def\Picardfunctor{\mathrm{Pic}}
\def\Deformationcategory{\mathcal{D}\!\mathit{ef}}
\usepackage{cleveref}
\crefname{originallemma}{original lemma}{original lemmas}
\Crefname{originallemma}{Original lemma}{Original lemmas}
\makeatletter
\newlabel{topologies-lemma-alternative}{{04HR}{source}{}{sourcepointer-1}{}}
\newlabel{topologies-lemma-alternative@cref}{{[originallemma][0][]04HR}{[1][0][]source}{}{}{}}
\makeatother
\crefname{originallemma}{original lemma}{original lemmas}
\Crefname{originallemma}{Original lemma}{Original lemmas}
\makeatletter
\newlabel{sites-lemma-colimit-sites}{{09YL}{source}{}{sourcepointer-2}{}}
\newlabel{sites-lemma-colimit-sites@cref}{{[originallemma][0][]09YL}{[1][0][]source}{}{}{}}
\makeatother
\crefname{originallemma}{original lemma}{original lemmas}
\Crefname{originallemma}{Original lemma}{Original lemmas}
\makeatletter
\newlabel{limits-lemma-descend-finite-presentation}{{01ZM}{source}{}{sourcepointer-3}{}}
\newlabel{limits-lemma-descend-finite-presentation@cref}{{[originallemma][0][]01ZM}{[1][0][]source}{}{}{}}
\makeatother
\crefname{originallemma}{original lemma}{original lemmas}
\Crefname{originallemma}{Original lemma}{Original lemmas}
\makeatletter
\newlabel{limits-lemma-limit-affine}{{01Z6}{source}{}{sourcepointer-4}{}}
\newlabel{limits-lemma-limit-affine@cref}{{[originallemma][0][]01Z6}{[1][0][]source}{}{}{}}
\makeatother
\crefname{originallemma}{original lemma}{original lemmas}
\Crefname{originallemma}{Original lemma}{Original lemmas}
\makeatletter
\newlabel{limits-lemma-descend-etale}{{07RP}{source}{}{sourcepointer-5}{}}
\newlabel{limits-lemma-descend-etale@cref}{{[originallemma][0][]07RP}{[1][0][]source}{}{}{}}
\makeatother
\crefname{originallemma}{original lemma}{original lemmas}
\Crefname{originallemma}{Original lemma}{Original lemmas}
\makeatletter
\newlabel{limits-lemma-descend-surjective}{{07RR}{source}{}{sourcepointer-6}{}}
\newlabel{limits-lemma-descend-surjective@cref}{{[originallemma][0][]07RR}{[1][0][]source}{}{}{}}
\makeatother
\crefname{originallemma}{original lemma}{original lemmas}
\Crefname{originallemma}{Original lemma}{Original lemmas}
\makeatletter
\newlabel{limits-lemma-descend-finite-presentation-variant}{{0EY1}{source}{}{sourcepointer-7}{}}
\newlabel{limits-lemma-descend-finite-presentation-variant@cref}{{[originallemma][0][]0EY1}{[1][0][]source}{}{}{}}
\makeatother
\title{The affine etale site of a directed limit of schemes is the colimit of the affine etale sites}
\author{The Stacks Project Authors}
\date{stacks-project:etale-cohomology.tex:local-snapshot-20261011; GNU FDL 1.2 or later}
\begin{document}
\maketitle
\noindent\textit{Editorial note.} All delivered mathematical text is the author's own, copied from the pinned Stacks Project chapter file etale-cohomology.tex; no statement or proof step was rewritten, reordered or removed.\par\medskip
\noindent\textit{Editorial note.} The proof is the author's complete argument for both cases, with the reduction from the general case to the quasi-compact one carried out rather than quoted.\par\medskip
\noindent\textit{Editorial note.} Every labelled reference occurring inside the counted spans is carried in assembly.external\_references as a non-load-bearing navigation pointer (same-chapter ones at their definition site, cross-chapter ones at the invocation inside the counted span), so no load-bearing mathematics is replaced by a stub. No bibliography citation occurs inside the counted spans.\par\medskip
\noindent\textit{Editorial note.} The checked chapter list for same-name collisions is the set of stacks chapters staged in this repository: chow, cohomology, divisors, etale-cohomology, flat, more-algebra, more-morphisms, morphisms, spaces-more-morphisms, varieties together with this card's chapter etale-cohomology.\par\medskip
\noindent\textit{Editorial note.} The preamble is the chapter's own native preamble excerpt already used for this source: the Xy-pic setup (native preamble lines 14-18) and the native theorem declarations with the mathematical macros (native preamble lines 154-217).\par\medskip
\noindent\textit{Allowed standard prerequisites (editorial).}\begin{itemize}
\item Affine etale site of a scheme. Topologies, the site whose objects are affine schemes etale over the scheme, with the etale topology. Supplies the two sides of the asserted equality of sites.
\item Colimit of sites along a directed inverse system. Sites, the construction of the colimit of a system of sites and the criterion for recognising an isomorphism of sites. Supplies the sense in which the displayed equality is asserted and the criterion used in the proof.
\item Affine transition morphisms and base change along a limit of schemes. Same chapter, the description of quasi-compact and quasi-separated objects of the affine etale site of a limit as base changes of objects over some index. Supplies the general case of the proof, in which an object of the limit site is descended to some stage.
\end{itemize}\medskip
\ifdefined\hypertarget\hypertarget{sourcepointer-1}{}\fi
\noindent\textit{Source reference.} Original lemma 04HR: sources/stacks-etale-cohomology/pinned-original-source.tex, source lines 7340--7340. Navigation pointer invoked at proof line 7340; the counted span contains the pointer only and the cited result is not reproduced inside it, so the pointer is registered as a non-load-bearing navigation stub in the established form for Stacks material. No bibliography citation is involved.\par\medskip
\ifdefined\hypertarget\hypertarget{sourcepointer-2}{}\fi
\noindent\textit{Source reference.} Original lemma 09YL: sources/stacks-etale-cohomology/pinned-original-source.tex, source lines 7345--7345. Navigation pointer invoked at proof line 7345; the counted span contains the pointer only and the cited result is not reproduced inside it, so the pointer is registered as a non-load-bearing navigation stub in the established form for Stacks material. No bibliography citation is involved.\par\medskip
\ifdefined\hypertarget\hypertarget{sourcepointer-3}{}\fi
\noindent\textit{Source reference.} Original lemma 01ZM: sources/stacks-etale-cohomology/pinned-original-source.tex, source lines 7357--7357. Navigation pointer invoked at proof line 7357; the counted span contains the pointer only and the cited result is not reproduced inside it, so the pointer is registered as a non-load-bearing navigation stub in the established form for Stacks material. No bibliography citation is involved.\par\medskip
\ifdefined\hypertarget\hypertarget{sourcepointer-4}{}\fi
\noindent\textit{Source reference.} Original lemma 01Z6: sources/stacks-etale-cohomology/pinned-original-source.tex, source lines 7360--7360. Navigation pointer invoked at proof line 7360; the counted span contains the pointer only and the cited result is not reproduced inside it, so the pointer is registered as a non-load-bearing navigation stub in the established form for Stacks material. No bibliography citation is involved.\par\medskip
\ifdefined\hypertarget\hypertarget{sourcepointer-5}{}\fi
\noindent\textit{Source reference.} Original lemma 07RP: sources/stacks-etale-cohomology/pinned-original-source.tex, source lines 7360--7360. Navigation pointer invoked at proof line 7360; the counted span contains the pointer only and the cited result is not reproduced inside it, so the pointer is registered as a non-load-bearing navigation stub in the established form for Stacks material. No bibliography citation is involved.\par\medskip
\ifdefined\hypertarget\hypertarget{sourcepointer-6}{}\fi
\noindent\textit{Source reference.} Original lemma 07RR: sources/stacks-etale-cohomology/pinned-original-source.tex, source lines 7366--7366. Navigation pointer invoked at proof line 7366; the counted span contains the pointer only and the cited result is not reproduced inside it, so the pointer is registered as a non-load-bearing navigation stub in the established form for Stacks material. No bibliography citation is involved.\par\medskip
\ifdefined\hypertarget\hypertarget{sourcepointer-7}{}\fi
\noindent\textit{Source reference.} Original lemma 0EY1: sources/stacks-etale-cohomology/pinned-original-source.tex, source lines 7373--7373. Navigation pointer invoked at proof line 7373; the counted span contains the pointer only and the cited result is not reproduced inside it, so the pointer is registered as a non-load-bearing navigation stub in the established form for Stacks material. No bibliography citation is involved.\par\medskip
\begin{lemma}
\label{lemma-colimit-affine-sites}
Let $I$ be a directed set. Let $(X_i, f_{i'i})$ be an inverse
system of schemes over $I$ with affine transition morphisms.
Let $X = \lim_{i \in I} X_i$. With
notation as in Topologies, Lemma \ref{topologies-lemma-alternative} we have
$$
X_{affine, \etale} = \colim (X_i)_{affine, \etale}
$$
as sites in the sense of
Sites, Lemma \ref{sites-lemma-colimit-sites}.
\end{lemma}

\begin{proof}
Let us first prove this when $X$ and $X_i$ are quasi-compact and
quasi-separated for all $i$ (as this is true in all cases of
interest). In this case any object of
$X_{affine, \etale}$, resp.\ $(X_i)_{affine, \etale}$ is
of finite presentation over $X$. Moreover, the
category of schemes of finite presentation over $X$ is the
colimit of the categories of schemes of finite presentation
over $X_i$, see
Limits, Lemma \ref{limits-lemma-descend-finite-presentation}.
The same holds for the subcategories of affine objects \'etale
over $X$ by Limits, Lemmas
\ref{limits-lemma-limit-affine} and \ref{limits-lemma-descend-etale}.
Finally, if $\{U^j \to U\}$ is a covering of $X_{affine, \etale}$
and if $U_i^j \to U_i$ is morphism of affine schemes \'etale over
$X_i$ whose base change to $X$ is $U^j \to U$, then we see that
the base change of $\{U^j_i \to U_i\}$ to some $X_{i'}$ is
a covering for $i'$ large enough, see
Limits, Lemma \ref{limits-lemma-descend-surjective}.

\medskip\noindent
In the general case, let $U$ be an object of $X_{affine, \etale}$.
Then $U \to X$ is \'etale and separated (as $U$ is separated)
but in general not quasi-compact. Still, $U \to X$ is locally
of finite presentation and hence by
Limits, Lemma \ref{limits-lemma-descend-finite-presentation-variant}
there exists an $i$, a quasi-compact and quasi-separated scheme $U_i$, and
a morphism $U_i \to X_i$ which is locally of finite presentation
whose base change to $X$ is $U \to X$. Then $U = \lim_{i' \geq i} U_{i'}$
where $U_{i'} = U_i \times_{X_i} X_{i'}$.
After increasing $i$ we may assume $U_i$ is affine, see
Limits, Lemma \ref{limits-lemma-limit-affine}.
To check that $U_i \to X_i$ is \'etale for $i$ sufficiently large,
choose a finite affine open covering $U_i = U_{i, 1} \cup \ldots \cup U_{i, m}$
such that $U_{i, j} \to U_i \to X_i$ maps into an affine open
$W_{i, j} \subset X_i$. Then we can apply
Limits, Lemma \ref{limits-lemma-descend-etale}
to see that $U_{i, j} \to W_{i, j}$ is \'etale
after possibly increasing $i$.
In this way we see that the functor
$\colim (X_i)_{affine, \etale} \to X_{affine, \etale}$
is essentially surjective. Fully faithfulness follows
directly from the already used
Limits, Lemma \ref{limits-lemma-descend-finite-presentation-variant}.
The statement on coverings is proved in exactly the same
manner as done in the first paragraph of the proof.
\end{proof}

\end{document}
